Propõe-se investigar a confiabilidade da inferência da massa do gráviton
e de polarizações adicionais de ondas gravitacionais com redes de
temporização de pulsares, conhecidas como PTAs. O trabalho dará
continuidade à análise de polarizações apresentada no TG de Wayne,
incorporando a descrição exata de ondas dispersivas, vínculos de teorias
de spin 2 massivo e métodos contemporâneos de inferência estatística. A
pergunta central será quanto a compressão da informação em frequência, a
escolha das distribuições a priori e as aproximações na resposta dos
pulsares modificam limites de massa e a capacidade de distinguir
polarizações.

O estudo partirá de resultados recentes sobre NANOGrav, MeerKAT e
perspectivas para SKA. Serão construídas simulações com parâmetros
conhecidos, comparando análises que preservam a frequência com análises
comprimidas. A validação incluirá reprodução de resultados analíticos,
convergência numérica, recuperação de sinais injetados, calibração
estatística e comparação entre informação anterior e posterior aos
dados. O resultado esperado é um conjunto reproduzível de critérios de
validade e de sensibilidade, com potencial para um artigo metodológico
em gravitação. Uma conclusão de que as aproximações são adequadas, ou de
que os dados ainda não distinguem os modelos, também será
cientificamente relevante se acompanhada de limites quantitativos.
